% !TEX root = ../nc-cubic.tex
\subsection{固定点計算による三重積の独立な検証}
\label{nft:sec:triple}
\begin{proposition}\label{nft:prop:triple}
定理\ref{nfc:thm:counterexample}の曲線 $T$ 上で
$\mathcal Q_b=\Ku(Y_b)^{G,\alpha}$ と書く。
$\rho=S^{-1}[3]$ とその自然変換空間 $R_j$ は同定理のものとする。
$\dim R_1(\mathcal Q_b)=7$、$\dim R_3(\mathcal Q_b)=35$ であり、
\[
 R_1(\mathcal Q_b)^{\otimes3}\longrightarrow R_3(\mathcal Q_b)
\]
の階数は $b\ne0$ で $35$、$b=0$ で $34$ である。
原点での余核は恒等元 sector の $[x_1x_2x_3]$ で生成される。
\end{proposition}
\begin{proof}
まず開集合を偏導関数から直接確認する。
共通零点では $w=0$, $x_i^2=-3u_i^2$ である。
一つでも $u_i=0$ なら、他の $x_j$ に関する偏導関数は
$2u_jx_j$ となるため、全座標が零となる。
それ以外では $c^2=-3$, $x_i=c\epsilon_i u_i$,
$\epsilon_i\in\{1,-1\}$、$e=\prod_i\epsilon_i$ と書く。
三つの残る式は $u_i^2=k u_j u_l$, $k=-bce/2$ であり、
積により $k^3=1$ が必要である。逆に $k^3=1$ なら
$u_1=1$, $u_2^3=1$, $u_3=k^{-1}/u_2$ と取れる。
従って射影特異点を持つパラメータは正確に $27b^6+64=0$ であり、
各パラメータでは符号積を指定した四組と $u_2$ の三通りから十二点を得る。
特に $0,2$ は同じ滑らかな係数開に属する。

次に $R_j$ を計算する。以下の計算には Clifford 演算子を使わない。
同変行列因子化の対角核の公式
\cite[Theorem~5.39 and its proof]{NFBFKKernels}を
$\mathbb G_m\times\widetilde G$ とその非自明な中心文字成分に適用する。
ここで $\widetilde G$ は生成元二つの持上げが反交換する
$G$ の二重被覆で、中心元を $z$ と書く。
曲率指標の核は $\mu_3\times\widetilde G$ である。
中央文字への射影の効果を明記する。誘導対角核のグラフ分解で
$z$ の作用は幾何的に同じグラフに付いた二つのラベル
$\widetilde a,z\widetilde a$ を交換する。
中心文字 $-1$ の冪等元 $(1-z)/2$ によりその差を取る。
従って $a\in G$ の sector を $v\in G$ で共役する際には、
二つの持上げの交換子による文字
$\chi_a(v)=\alpha(a,v)/\alpha(v,a)$ が付く。
これはねじれ群環の基底で $e_v e_a e_v^{-1}$ を取ることでも確かめられる。

$\xi\in\mu_3$ が非自明なら $\xi a$ の移動空間の次元は七である。
固定部分の Jacobian Koszul 複体は次数零にのみコホモロジーを持つので、
この奇余次元の sector はコホモロジー次数零の $R_j$ に寄与しない。
$\xi=1$, $a=1$ の項は $\operatorname{Jac}(F_b)_j^G$ である。
$a\ne1$ は二つの座標を反転する。そこに残る active 座標を $x_k$ とすると
固定式は $F^{a}=u_k x_k^2+g$ であり、臨界点は原点だけである。
余次元二の determinant と三次周期の次数補正は $j+2-3=j-1$ となる。
法線二座標の行列式指標は、核が $\langle a\rangle$ となる唯一の非自明指標であり、$\chi_a$ と相殺する。
以上から
\[
 R_j(\mathcal Q_b)=
 \operatorname{Jac}(F_b)_j^G\ \oplus \bigoplus_{1\ne a\in G}\operatorname{Jac}(F^{a})_{j-1}^G
\]
というベクトル空間分解を得る。恒等項にある不変多項式の乗法は各項で
その多項式の固定空間への制限の乗法である。実際、対角とグラフの
導来交差を計算する同公式の Koszul 複体では、その多項式は係数として作用し、
移動座標を消す射は固定空間への制限である。他の二つの sector 同士の積の式は不要である。

次数一では恒等項に $U=(u_1,u_2,u_3,w)$ の四座標があり、
三つの非恒等項には定数が一つずつある。
次数三の恒等項では、二次関係
$x_i^2=-3u_i^2$, $w^2=0$ により、$U$ の三次単項式十六個と
$P=x_1x_2x_3$ だけが残る。残りの関係は三つの
\[
                     6u_i^3=bP .
\]
実際、勾配イデアルの次数三において、$\partial_{x_i} F_b$ に
掛けられる逆指標の一次式は $x_i$ の定数倍に限られ、
$\partial_{u_i} F_b,\partial_w F_b$ には $U$ の一次式だけが掛かる。
これで全不変関係を数えており、表示した三関係は独立である。
従って恒等項は十四次元であり、$U$ だけの三重積の像は
$b\ne0$ では十四次元、$b=0$ では十三次元となる。

非恒等項の次数三では $\operatorname{Jac}(F^a)_2^G$ を取る。
$x_k^2=-3u_k^2$、$u_j^2=0$ ($j\ne k$)、$w^2=0$ より、
$U$ の二次式の七次元商を得る。$u_k x_k=0$ は非自明指標の関係なので
この次数の不変部分に追加の関係を与えない。
したがって各項は対応する $R_1$ の定数と $U$ の二元の積で張られる。
以上により $\dim R_3=14+3\cdot7=35$ である。
$b\ne0$ では既に全像、原点でも少なくとも $34$ 次元の像を得た。

原点では $x_1,x_2,x_3$ を同時に負にし $U$ を固定する対合 $h$ が
$F_0$ を保つ。$h$ は群作用と可換し、中心拡大には自明に、
次数シフトには自然に作用する。恒等項の $R_1$ は $h$ 不変である。
非恒等三項でも定数が残り、$h$ は法線二方向をともに負にするので
その行列式は一となる。従って $h$ は $R_1$ 全体に自明に作用し、
いかなる三重積も $h$ 不変である。他方 $P$ は $h$ で符号を変え、
原点の恒等項で非零である。例えば $U=0$ に特殊化した
$\C[x_1,x_2,x_3]/(x_1^2,x_2^2,x_3^2)$ でも非零である。
これで階数と余核を得る。

任意の滑らかな可換三次五様体では、七変数に対する同じ公式で
非恒等 scalar sector が奇余次元のため消え、$R_1,R_3$ と積は通常の
Jacobian 環のものになる
\cite[Theorem~3.5, Remark~3.6, Proposition~3.7]{LZ}。
この環は一次生成なので三重積は全射である。
よって原点の非実現は、定理\ref{nfc:thm:counterexample}の六つの
二次関係の具体的な構成から独立にも従う。対合を標識として比較先へ
運ぶ必要はなく、三重積の階数は $\rho=S^{-1}[3]$ から決まる。
\end{proof}
